\section{The mechanism}
\label{sec:mech}
\begin{table*}[t]
\centering\small
\caption{On-chain invariants checked on every rebalancing fill.}
\label{tab:inv}
\begin{tabular}{L{22mm}L{56mm}L{44mm}L{30mm}}
\toprule
Rule & Check & Stops & Default\\
\midrule
Drift must fall & $D_{\mathrm{after}}\le D_{\mathrm{before}}-\Delta_{\min}$; no token overshoots its target by more than $\varepsilon$ & Fills that push the basket away from target & $\Delta_{\min}=0.25\%$ of NAV\\
Price bound & Value in $\ge(1-s_{\max})\times$ value out, both at oracle prices & Overcharging & $s_{\max}=30$\bps\\
Oracle freshness & One price set per fill, used for both trigger and price bound; age $\le\tau_{\max}$; timestamps not older than the newest the vault has used; feeds within $\theta$ & Fills at stale, selected or manipulated prices (Section~\ref{sec:oracle}) & $\tau_{\max}=2$\,s; $\theta=50$\bps\\
Size caps & Fill $\le q_{\max}\times$NAV; daily total $\le Q_{\max}\times$NAV & Draining the vault & 5\% per fill, 20\% per day\\
\bottomrule
\end{tabular}
\end{table*}

Section~\ref{sec:theory} says what a gated pool must do: move its inventory only when drift reaches a band that is chosen, not implied by a fee, and pay a competitive spread when it does. This section specifies a design that does that and bounds what each privileged party can take. Figure~\ref{fig:arch} shows one pool.

\begin{figure}[t]
\centering
% Figure 1: architecture of one gated pool (black and white, TikZ)
\begin{tikzpicture}[x=1mm,y=1mm,
  box/.style={draw,line width=0.5pt,minimum width=22.6mm,minimum height=11mm,align=center,inner sep=1.2pt,font=\scriptsize,fill=white},
  key/.style={box,double,double distance=0.9pt},
  ttl/.style={font=\scriptsize\bfseries},
  lab/.style={font=\tiny\itshape,inner sep=1.2pt},
  tip/.style={-{Stealth[length=1.7mm,width=1.2mm]},line width=0.45pt},
  both/.style={{Stealth[length=1.7mm,width=1.2mm]}-{Stealth[length=1.7mm,width=1.2mm]},line width=0.45pt}]
\node[box] (ret) at (0,0)   {\textbf{Retail traders}\\anyone, size-capped\\pay fee and impact};
\node[key] (quo) at (29.2,0)  {\textbf{Quote layer}\\thin weighted curve\\centred on the oracle};
\node[box] (ora) at (58.4,0)  {\textbf{Oracle adapter}\\two price feeds\\stale or split: halt};
\node[box] (lps) at (0,-19) {\textbf{Liquidity providers}\\deposit and redeem\\in kind, any time};
\node[box] (vau) at (29.2,-19){\textbf{Vault inventory}\\basket at weights $w$\\keeps fee and impact};
\node[box] (sol) at (29.2,-38){\textbf{Bonded solvers}\\Dutch auction\\fill at oracle $\pm s$};
\node[box] (ext) at (58.4,-38){\textbf{External venue}\\exchange or perps\\hedges each fill};
\draw[both] (ret) -- node[lab,above]{swaps} (quo);
\draw[tip]  (ora) -- node[lab,above]{anchor} (quo);
\draw[both] (lps) -- node[lab,above]{in kind} (vau);
\draw[both] (sol) -- node[lab,above]{hedge} (ext);
\draw[both] (quo) -- node[lab,right,align=left]{re-centred every\\block, internally} (vau);
\draw[both] (vau) -- node[lab,left,align=right]{only when\\$D>\delta$} (sol);
\draw[tip] (ora.south) -- (58.4,-28.5) node[lab,left,align=right,pos=0.62]{bounds every\\fill at $\pm s$} -- (37,-28.5) -- (37,-32.5);
\end{tikzpicture}

\caption{One gated pool. Quotes follow the oracle every block through internal re-centring; inventory moves only when drift exceeds the band $\delta$, through bonded solvers at oracle price $\pm s$.}
\label{fig:arch}
\end{figure}

\subsection{Vault, NAV and drift}
A vault holds $x_i$ of each token. It values itself at oracle prices $p_i$ and measures its drift from target:
\begin{equation}
\NAV=\sum_i x_i p_i,\qquad \omega_i=\frac{x_i p_i}{\NAV},\qquad D=\tfrac12\sum_i\lvert\omega_i-w_i\rvert .
\label{eq:drift}
\end{equation}
$D$ is the one-way notional, as a share of NAV, that restores target weights. LPs mint and redeem in kind: depositing a fraction $\alpha$ of every balance mints $\alpha$ of the share supply, and burning shares returns the same fraction of every balance. Neither path needs a price, so neither can be gamed through the oracle or paused.

\subsection{The gate}
\label{sec:gate}
\paragraph{Trigger.} Every interaction recomputes $D$, and anyone may call a public check. When $D>\delta$ the vault opens a rebalancing intent listing the token deltas that move drift to $\din$.

\paragraph{Reverse Dutch auction.} The spread a solver may keep starts at zero and rises with time, $s(t)=\min(s_{\max},\,\varrho\, t)$, with $\varrho$ = 1\bps\ per second. The first whitelisted solver whose cost is covered fills. With several solvers the fill lands near the cheapest solver's cost plus one step; with one solver it lands at $s_{\max}$, which Section~\ref{sec:costs} treats as the stress case. A solver's cost is the hedge, the same $c$ an arbitrageur pays on a public pool, plus inventory risk on a larger fill.

\paragraph{On-chain invariants.} Every fill must pass four checks, whoever submits it (Table~\ref{tab:inv}). Together they bound the damage of a compromised or hostile solver to $s_{\max}Q_{\max}=0.3\%\times20\%=0.06\%$ of NAV a day, plus any oracle error (Appendix~\ref{app:ex}, B13).


\paragraph{Bonds and fallback.} Solvers post USDC bonds sized to seven days of worst-case damage, so a \$100k bond allows $B/(7s_{\max})\approx\$4.76$M of fills a day. Bonds are slashed only on objective on-chain evidence. If no whitelisted solver fills within $T_{\mathrm{fb}}$ (one hour), the same auction opens to anyone who posts a bond, so the pool stays live without trusting any operator, and a colluding whitelist cannot charge more than outsiders would.

\subsection{The quote layer}
\label{sec:quote}
Closing the inventory to arbitrage would also close it to retail. The quote layer reopens it without reopening arbitrage.

\paragraph{Anchored curve.} At the start of each block the layer sets virtual reserves $\hat x_i=w_i\NAV/p_i$ centred on the oracle. A retail trade executes on these reserves with weighted-pool math and fee $f_q$, and the vault's real balances change by the same amounts; trades in one block move along the curve, and the next block starts again from the oracle. Re-centring is an internal transfer, so the quote stays fresh without the inventory trading. In a public pool the next arbitrageur reverses a retail trade's price impact and keeps it; here the impact stays with LPs, and the inventory change it causes is netted by later flow or rebalanced at the band. A small purchase of token $i$ with token $n$ moves the quote by
\begin{equation}
\Delta\ln P_i\approx\frac{A}{\NAV}\Big(\frac{1}{w_n}+\frac{1}{w_i}\Big),
\label{eq:impact}
\end{equation}
so \$10k of HYPE (10\% weight) paid in USDC (15\%) in a \$10M vault moves it 1.67\% and pays \$30 of fee and \$82 of impact (B10). Because re-centring would let a large trader split an order across blocks, the anchor carries an \emph{inventory skew}: the centre for token $i$ is shifted against the side the vault is already short of, in proportion to $(w_i-\omega_i)/w_i$.

\paragraph{What the quote layer cannot do.} It quotes at a price someone else measured, so it is exposed to anyone who knows the price better or sooner. Section~\ref{sec:oracle} quantifies this. The layer therefore enforces a freshness gate (no trades while the price is older than $\tau_{\max}$ or the feeds disagree), monotone oracle timestamps, a per-window size cap (default 0.1\% of TVL per 12 seconds), and a fee floor tied to realised volatility. The strongest option is \emph{next-price execution}: a trade is submitted in one block and filled at the first oracle price stamped after it, which removes latency and selection edges at the cost of a one-block delay.
